Under CECL a lender may estimate its allowance for credit losses with any
reasonable method, and the SEC expects it to document an identified range of
expected losses. This paper measures how much of that range comes from the
choice among equally accurate default models. On \num{data.kept} Freddie
Mac mortgages, at each of \num{const.n_dates} year-ends from 2005 to 2025, I
train \num{const.n_models} models on what a lender knew at the time, keep those
within one bootstrap standard error of the best held-out log loss, and compute
the lifetime allowance of the same portfolio under each, with the prepayment
model, the economic forecast, the reversion policy and loss severity held
fixed. The design and six predictions were recorded before the data were
opened. Across these models the allowance ranges over a median of
\num{res.rel.median}\% of the most accurate model's (\num{ph.paired2.rel}\% and
\num{res.paired1.rel}\% among the few models that a paired test at two or one
standard errors cannot separate from the best), up to \num{res.usd.max} billion dollars on
Freddie Mac's active book, and more than the gap between one- and three-year
reversion at every date. In comparisons added after the results, refitting the best
model on independent samples of the same size moves the allowance at least as
far as switching among twenty equally accurate models (\num{p2.disj.rsm}\%
against \num{p2.disj.lfl}\% where the samples are disjoint) and changes more
loans' IFRS 9 stage, contrary to the interpretation of a staging prediction
that holds as recorded. Realized losses fall inside the
range at one of \num{H5.n_elig} dates: model uncertainty is not forecast
uncertainty. The spread was not unusually large in relative terms before the
crisis, and monotone constraints barely narrow it. Code and every aggregate
result are released.
